% Effective perfect C1-field extension; full written audits by root and dieudonne passed.
\subsection{C.3 over perfect \texorpdfstring{$C_1$}{C1} fields with explicit arithmetic access}

A field $F$ is $C_1$ if every homogeneous polynomial of degree $d>0$
in more than $d$ variables has a nonzero zero over $F$.
We first give the field-theoretic classification, and then specify the
extra arithmetic access under which every operation is a terminating
algorithm. These are separate assertions.

\begin{lemma}[Known $C_1$ form arithmetic, rechecked]
\label{C:C1_field_forms}
Let $F$ be a $C_1$ field of characteristic different from two.
Every finite extension of $F$ is $C_1$ and has trivial Brauer group.
For every such extension $K$, nonsingular symmetric forms over $K$
are classified by their dimension and determinant square class.
Every quadratic field extension $K/K_0$ of finite extensions of $F$
has surjective norm, and nonsingular hermitian forms for its nontrivial
involution are classified by dimension. Consequently
\[
 \begin{split}
 GW^+(K,1)&\simeq\mathbb Z\oplus K^\times/(K^\times)^2,\\
 GW^-(K,1)&\simeq\mathbb Z,\qquad
 GW^\varepsilon(K,\sigma)\simeq\mathbb Z\quad(\sigma\ne1).
 \end{split}                                               \tag{C.3.Q1}
\]
The first isomorphism uses rank and the ordinary determinant; the
square-class factor is written multiplicatively. The second generator
is a symplectic plane, and the last generator has multiplicity dimension
one, after the usual scalar change for a skew-hermitian form.
\end{lemma}
\begin{proof}
If $L/F$ has degree $a$ and $f$ is homogeneous of degree $d<n$ over $L$,
expand its $n$ variables in an $F$-basis of $L$. The polynomial
$N_{L/F}(f)$ is homogeneous of degree $ad$ in $an$ variables over $F$.
A nonzero zero gives a nonzero zero of $f$, since the norm of a nonzero
field element is nonzero. Thus $L$ is $C_1$. The reduced-norm proof of
\cite[Theorem 59.67.8, Tag 03RC]{C-StacksC1Fields} gives
$\operatorname{Br}(L)=0$: a central division algebra of degree $b>1$
would have a reduced norm of degree $b$ in $b^2$ variables with no
nonzero zero, contradicting the $C_1$ property.

For a quadratic extension $K/K_0$ and $c\in K_0^\times$, the ternary
quadratic equation $N_{K/K_0}(x)-cz^2=0$ has a nonzero solution.
Such a solution has $z\ne0$, because $N(x)=0$ implies $x=0$.
It therefore gives $N(x/z)=c$. This proves norm surjectivity.

Every nonsingular binary quadratic form is a scalar multiple of a
quadratic-field norm form, or is hyperbolic. In the first case norm
surjectivity absorbs its scalar; in the second case the hyperbolic
plane has one isometry class. Hence binary forms with the same
determinant square class are isometric. All forms of dimension at
least three are isotropic by $C_1$, and splitting a hyperbolic plane
then proves the rank--determinant classification by induction.
Dimension one is the square-class classification itself.
Orthogonal sum adds ranks and multiplies determinants. The group
completion is exactly the first group in (C.3.Q1): the rank-zero
coordinate of $c$ is represented by
$[\langle c\rangle]-[\langle1\rangle]$.

Alternating forms have a symplectic basis. For a nontrivial field
involution, hermitian Gram--Schmidt diagonalizes a nonsingular form
with diagonal entries in the fixed field. Norm surjectivity normalizes
each diagonal entry to one. A nonzero element $c$ with
$\sigma(c)=-c$ changes skew-hermitian to hermitian forms; such an
element is $b-\sigma(b)$ for any $b$ not fixed by $\sigma$.
These facts prove the remaining assertions.
\end{proof}

\begin{proposition}[Proved: the arithmetic form parameters over every perfect $C_1$ field]
\label{C:C1_Galois_WG}
Let $k$ be perfect and $C_1$ of characteristic $p>2$.
In Theorem~\ref{C:perfect_Galois_indices} all Schur indices are one.
For each self-dual Galois index $i$ choose an invariant nonsingular
$\varepsilon_i$-symmetric form $q_i$ on its associated $k$-module $M_i$.
Its residue endomorphism algebra is a finite field extension $K_i/k$;
let $\sigma_i$ be the adjoint involution on that field. Partition
the indices into symmetric type $(\sigma_i=1,\varepsilon_i=1)$,
alternating type $(\sigma_i=1,\varepsilon_i=-1)$, unitary type
$(\sigma_i\ne1)$, and one index from every non-self-dual pair.
Then the split orthogonal-sum group is
\[
 WG_{\rm sp}(BO_m/k)\simeq
 \bigoplus_{i\in\mathscr J_{\rm s}}
    \bigl(\mathbb Z\oplus K_i^\times/(K_i^\times)^2\bigr)
 \ \oplus\!
 \bigoplus_{i\in\mathscr J_{\rm a}\cup\mathscr J_{\rm u}
                      \cup\mathscr J_{\rm h}}\mathbb Z .    \tag{C.3.Q2}
\]
The integer in symmetric or unitary type is multiplicity dimension;
in alternating type it is half that dimension; in a paired type it
is the common multiplicity of the two members. In symmetric type
the second coordinate is the determinant of the residual multiplicity
form. These are explicit field-form parameters, without a residual
hermitian-form isometry quotient.
\end{proposition}
\begin{proof}
The center of each residue division algebra is finite over $k$.
Lemma~\ref{C:C1_field_forms} makes that division algebra equal to its
center. Equation (C.3.G2) therefore gives Schur index one.
The checked hypotheses of the hermitian-category theorem
\cite[Propositions 2.4--2.5 and Theorems 3.2--3.3,
pp.271--273]{C-QSS1979} hold for finite-dimensional rational modules:
idempotents split, endomorphism rings are finite dimensional, and their
radicals are nilpotent. Transfer and radical reduction give the
multiplicity-form groups over $K_i$. The fixed field of a nontrivial
$\sigma_i$ is a finite extension of $k$ and is $C_1$ by the lemma.
Thus (C.3.Q1) gives (C.3.Q2), including the non-self-dual hyperbolic
integer factors. The Galois indices and finite-field data are those
of Theorem~\ref{C:perfect_Galois_indices}; no descent of an absolute
indecomposable to a field of moduli is asserted without that theorem.
\end{proof}

\paragraph{Effective field data.}
For the algorithmic statements below, $k$ is an \emph{infinite countable}
perfect $C_1$ field of characteristic $p>2$, supplied with a computable
surjective enumeration and exact, decidable field arithmetic and equality.
We also assume univariate factorization over $k$ and over each explicitly
presented finite extension $k[z]/(f)$ is available. A repeated enumeration
can be made injective by retaining an element only at its first occurrence;
we use that resulting enumeration as the order on coefficients. These
are arithmetic hypotheses, not conclusions from the abstract $C_1$
property. A concrete infinite nonclosed family with these procedures
is given below.

\begin{lemma}[Proved: terminating certificates for finite endomorphism algebras]
\label{C:C1_algebra_certificates}
With the effective field data just stated, an algorithm computes the
Jacobson radical of every finite-dimensional unital $k$-algebra $E$
given by a multiplication table, together with an explicit isomorphism
\[
 E/\operatorname{rad}E\ \simeq\
          \prod_{j=1}^t M_{r_j}\bigl(k[z]/(f_j)\bigr),      \tag{C.3.Q3}
\]
where the $f_j$ are monic irreducible polynomials. It also computes
primitive idempotent lifts. In particular, it decides whether $E$
is local. When $E$ is the endomorphism algebra of a given rational
representation, it produces its indecomposable summands.
\end{lemma}
\begin{proof}
Put $h=\dim_k E$, and handle the zero algebra separately. Enumerate
all finite lists of positive integers $r_j,d_j$ with
$\sum_jr_j^2d_j\le h$, all monic polynomials $f_j$ of degree $d_j$,
and all $k$-linear maps from $E$ to the corresponding displayed product.
This is a countable effective enumeration by the field access above.
Factorization decides whether each $f_j$ is irreducible. Retain a
candidate map $\pi$ exactly if it is unital, multiplicative on a basis,
surjective, and its kernel $J$ satisfies $J^{h+1}=0$.
These are finite matrix, rank and ideal-product tests.

Every retained map has $J=\operatorname{rad}E$: its kernel is a
nilpotent ideal and its quotient is semisimple. Such a map exists.
Indeed, Wedderburn decomposition gives the semisimple quotient as
matrix algebras over finite-dimensional division algebras. Each center
is a finite extension of $k$ and is $C_1$, so the division algebra
is its center. Perfectness makes these centers separable, hence
primitive and of the form $k[z]/(f_j)$. The nilpotent radical has
nilpotency index at most $h+1$. Therefore the enumeration eventually
finds a verified map, proving termination; no negative search for an
idempotent or a rational point is used.

For a standard primitive matrix idempotent $\bar e$ in the target,
choose a preimage $e\in E$. Put $N=h+1$. Then
$(e^2-e)^N=0$. The extended Euclidean algorithm supplies a polynomial
$q(T)$ congruent to zero modulo $T^N$ and to one modulo $(T-1)^N$.
Consequently $q(e)^2=q(e)$ and $\pi(q(e))=\bar e$.
Its corner quotient is a field, so this lift is primitive.
When $E$ is the representation's equivariant endomorphism algebra,
split into the invariant image and kernel of this matrix and continue
on the complementary summand; positive dimensions
decrease. This gives all required indecomposable summands in finitely
many such steps. Finally $E$ is local precisely when the quotient
in (C.3.Q3) is one field factor with $r_1=1$.
\end{proof}

\begin{theorem}[Proved: complete recursive indices and operations over effective perfect $C_1$ fields]
\label{C:Om_effective_C1}
For every field with the effective data above and every split $O_m/k$,
there is a decidable, repetition-free set of matrix indices for all
indecomposable rational representations. An index specifies dimension,
minimal coefficient degree, and a canonical natural-number code for
its coefficient matrices. Duality, all data in (C.3.Q2), equality of
form coordinates, all products, and every prescribed exterior-power
coefficient of a finite virtual input are computable in finite time.

The square-class coordinates in (C.3.Q2) have canonical representatives
in the explicit enumeration of $K_i$. Thus no undecided representation
or form isomorphism quotient remains. The indices at fixed dimension
and degree may be infinite; this is a recursive matrix classification,
not an assertion of finitely many polynomial templates over $k$.
\end{theorem}
\begin{proof}
\emph{Decidable module indices.}
Construct the standard-monomial coalgebras $C_d\subset k[O_m]$ and their
finite dual-algebra multiplication constants as in
Theorem~\ref{C:Om-canonical-templates}. For fixed $n,d$, tuples of
$r=\dim C_d$ matrices in $M_n(k)$ have a computable one-to-one coding
by natural numbers, using the coefficient enumeration. The multiplication,
unit and minimal-degree tests (C.3.T3)--(C.3.T4) are decidable equalities.
The centralizer is a finite-dimensional algebra computed by linear
equations. Lemma~\ref{C:C1_algebra_certificates} decides whether it is
local, hence whether the module is indecomposable, and constructs a
decomposition when needed.

Two valid tuples of dimension $n$ are isomorphic precisely when the
determinant polynomial on their linear intertwiner space is nonzero.
Its coefficients are computed by finite symbolic polynomial arithmetic.
If nonzero, a grid of $n+1$ distinct enumerated field elements gives an
invertible intertwiner, as in Lemma~\ref{C:NoetherDeuring_rational}.
For a tuple with code $c$, test the finitely many codes $0,\ldots,c$
and retain the least isomorphic one. This is the least code in the
entire orbit, since the input code is itself in that orbit. Thus the
canonical condition is decidable and gives exactly one index per
isomorphism class. Degree is basis invariant; the minimal-degree test
removes repetition across $d$. Every rational representation occurs
at a finite degree.

\emph{Duality and the field-form parameters.}
The antipode gives the dual tuple, which is canonically indexed by the
same finite procedure. This decides self-duality and forms the paired
indices. For a self-dual module $M_i$, solve the linear invariant-form
equations in the symmetric and alternating matrix spaces. The determinant
polynomial decides whether either space contains a nonsingular form.
Choose symmetric first if possible; otherwise choose alternating. QSS
guarantees success. A finite grid gives the chosen actual matrix $q_i$.

Apply the certificate lemma to $E_i=\operatorname{End}(M_i)$. Since
$M_i$ is indecomposable, its explicit quotient is a field
$K_i=k[z]/(f_i)$. Compute the adjoint involution on $E_i$ using $q_i$,
and its induced map $\sigma_i$ on $K_i$ by linear preimages and projection.
This determines all four types in (C.3.Q2) by finite tests.

Square-class equality in $K_i$ is decidable by factoring $X^2-a/b$.
For an input nonzero $a$, scan the finitely many nonzero enumeration
codes no larger than its own and retain the least representative of
its square class. This gives a canonical scalar for every square-class
coordinate. Products and inverses of these coordinates are ordinary
field operations followed by that finite normalization.

\emph{Actual representatives for group coordinates.}
In symmetric type the integer generator is $[q_i]$. For
$a\in K_i^\times$, choose a lift $b\in E_i$, replace it by
$a'=(b+b^*)/2$, and use
$[q_i a']-[q_i]$ for the rank-zero square-class coordinate of $a$.
Its reduction is $a$, so $a'$ is invertible and the form is symmetric.
For alternating and paired types use $H(M_i)$. In unitary type choose
$0\ne c\in K_i$ with $\sigma_i(c)=\varepsilon_i c$, lift it to $b$,
and set $a'=(b+\varepsilon_i b^*)/2$. Then $q_i a'$ is symmetric and
nonsingular and represents the integer generator. For the negative
sign, a nonzero difference $c=b_0-\sigma_i(b_0)$ is found among a
$k$-basis of $K_i$; for the positive sign use $c=1$.
Every finite coordinate input is consequently an explicit finite
integral combination of actual forms.

\emph{Extracting coordinates, not searching an isometry quotient.}
For an actual equivariant symmetric form, decompose its underlying
representation by the certificate lemma, canonically index the summands,
and construct the isomorphisms to the chosen representatives by the
intertwiner grid. Group equal indices into isotypic blocks. For a
self-dual block $M_i^{\oplus b}$, let $B_{uv}:M_i\to M_i^*$ be
the corresponding blocks of its Gram matrix. The entries
\[
        H_{uv}=\pi_i(q_i^{-1}B_{uv})\in K_i              \tag{C.3.Q4}
\]
form a nonsingular $\varepsilon_i$-hermitian matrix for $\sigma_i$.
To see nonsingularity, reduce the global isomorphism to the semisimple
quotient of the endomorphism category. Nonisomorphic indecomposable
blocks have only radical off-diagonal morphisms; the duality fixes this
self-dual block. Its quotient block must therefore be invertible.
The radical-reduction theorem of QSS identifies (C.3.Q4) with its
multiplicity form. In symmetric type its coordinates are
$(b,\det H)$, with the determinant normalized modulo squares; in
alternating type $b$ is even and the coordinate is $b/2$; in unitary
type the coordinate is $b$. In a non-self-dual pair the two
multiplicities agree and give the hyperbolic coordinate. This computes
the complete class using finite matrices and the already proved
field-form classification.

Finally tensor products multiply coaction entries and take tensor Gram
matrices; exterior powers take coefficient minors and determinant Gram
forms. Their degree bounds are $d+e$ and $jd$, respectively, and no
factorial is inverted. Apply the preceding extraction to each resulting
actual form. Bilinearity and coefficientwise inversion of
$\lambda_t(-x)=\lambda_t(x)^{-1}$ compute virtual products and any
specified finite set of exterior coefficients. All subroutines either
are finite operations or are positive certificate searches whose
termination has been proved. This proves the theorem.
\end{proof}

\begin{corollary}[Proved: a concrete infinite nonclosed effective family]
\label{C:C1_tower_example}
Fix an odd prime $p$ and any prime $\ell$. In a fixed algebraic closure
of $\mathbb F_p$ put
\[
                 k_{p,\ell}=\bigcup_{a\ge0}\mathbb F_{p^{\ell^a}}.
                                                               \tag{C.3.Q5}
\]
This infinite perfect nonclosed field has a uniform computable
presentation with all the effective data required above, including
factorization over every specified finite extension. Thus
Theorem~\ref{C:Om_effective_C1} gives an ordinary terminating full
split-$WG$ algorithm over each such fixed field.
\end{corollary}
\begin{proof}
The finite-field tower and coherent embeddings are computable by finite
search, as in Theorem~\ref{C:locally-finite-equality-obstruction}.
Stage-element names give decidable arithmetic and equality, and the
first-occurrence procedure gives an injective enumeration. The field
is perfect, and it is not algebraically closed because it does not
contain $\mathbb F_{p^r}$ for any prime $r\ne\ell$.

Finite fields are $C_1$. For completeness, if a homogeneous polynomial
of degree $d<n$ over $\mathbb F_q$ had only the zero zero, sum
$1-f(x)^{q-1}$ over $x\in\mathbb F_q^n$. Every monomial of
$f^{q-1}$ has total degree less than $n(q-1)$, hence has one exponent
less than $q-1$; summation in that variable is zero, including exponent
zero since $q=0$ in the field. The resulting number of zeros is zero
modulo $p$, whereas a unique zero would give one. This contradiction
proves $C_1$. Coefficients of any polynomial over (C.3.Q5) lie in a
finite stage, so the union is $C_1$ as well.

To factor a polynomial, place its coefficients in one finite stage and
factor there. An irreducible factor of degree $d$ over that stage splits
under an extension of degree $\ell^b$ into $\gcd(d,\ell^b)$ factors,
each of degree $d/\gcd(d,\ell^b)$. This is the Frobenius-orbit formula
for finite fields. Taking $b$ at least the largest $\ell$-adic valuation
of these finitely many degrees makes all remaining factor degrees
prime to $\ell$, so no further tower stage can split them. The resulting
factorization is therefore final over $k_{p,\ell}$.

Any specified finite extension is algebraic over $\mathbb F_p$ and
generated over $k_{p,\ell}$ by one element lying in a finite field.
Its presentation embeds effectively in a tower
$\bigcup_a\mathbb F_{p^{D\ell^a}}$: put its defining coefficients
in a finite field, choose a root in a finite extension by finite search,
and use the resulting finite-field arithmetic and the injectivity of
the given irreducible presentation. The same degree argument gives
factorization over this tower. This supplies the finite-extension
arithmetic required in the theorem. No undecidable future change of
tower degrees occurs in (C.3.Q5).
\end{proof}

\paragraph{Further structural examples and scope.}
By Tsen's theorem \cite[Theorem 59.67.10, Tag 03RD]{C-StacksC1Fields},
the function field of a curve over an algebraically closed field is
$C_1$. Its algebraic extensions are $C_1$ by the finite-extension
norm proof and finite coefficient descent; in particular its perfect
closure has the structural formula (C.3.Q2). The ordinary algorithm is
asserted only when the extra effective arithmetic data are supplied.
It does not contradict Theorem~\ref{C:locally-finite-equality-obstruction},
whose presentations do not uniformly supply root or factorization
decisions. General perfect fields with nontrivial residue division
algebras, and general nonfield bases, retain their separate gates.
